% ── Shape notation (per CLAUDE.md "Shape notation — always tag the algebra") ──
%
% K⟨n,m,p⟩          non-commutative matmul over K
% K⟨n,m,p⟩ᶜ         commutative matmul over K
% K⟨n,m,p⟩:m        multiplications count
% K⟨n,m,p⟩:r        rank
% K⟨n,m,p⟩:a        addition count
%
% Implemented as plain LaTeX commands so they render cleanly in the
% paper. Authors should USE these macros and not bare angle brackets.

\newcommand{\nmpshape}[3]{\ensuremath{\langle #1, #2, #3 \rangle}}
\newcommand{\nmpshapec}[3]{\ensuremath{\langle #1, #2, #3 \rangle^{\mathrm{c}}}}% field-less commutative shape
\newcommand{\nmpfield}[4]{\ensuremath{#1\langle #2, #3, #4 \rangle}}     % K⟨n,m,p⟩
\newcommand{\nmpfieldc}[4]{\ensuremath{#1\langle #2, #3, #4 \rangle^{\mathrm{c}}}}% K⟨n,m,p⟩ᶜ
\newcommand{\nmpcount}[5]{\ensuremath{#1\langle #2, #3, #4 \rangle\!:\!#5}} % suffix :m / :r / :a

% Common short forms
\newcommand{\Ff}{\ensuremath{\mathbb{F}}}
\newcommand{\Ftwo}{\ensuremath{\Ff_2}}
\newcommand{\Fthree}{\ensuremath{\Ff_3}}
\newcommand{\Reals}{\ensuremath{\mathbb{R}}}
\newcommand{\Complex}{\ensuremath{\mathbb{C}}}
\newcommand{\Rationals}{\ensuremath{\mathbb{Q}}}
\newcommand{\Integers}{\ensuremath{\mathbb{Z}}}

% Theorem environments
\newtheorem{theorem}{Theorem}
\newtheorem{lemma}[theorem]{Lemma}
\newtheorem{proposition}[theorem]{Proposition}
\newtheorem{corollary}[theorem]{Corollary}
\theoremstyle{definition}
\newtheorem{definition}[theorem]{Definition}
\newtheorem{example}[theorem]{Example}
\theoremstyle{remark}
\newtheorem{remark}[theorem]{Remark}

% Style for code-ish identifiers (file names, JSON keys, Java symbols)
\newcommand{\code}[1]{\texttt{#1}}
